
\documentclass{article}
\usepackage{colortbl}
\usepackage{makecell}
\usepackage{multirow}
\usepackage{supertabular}

\begin{document}

\newcounter{utterance}

\centering \large Interaction Transcript for game `hot\_air\_balloon', experiment `air\_balloon\_survival\_bg\_complexity\_hard', episode 3 with mistral{-}large{-}3{-}azure.
\vspace{24pt}

{ \footnotesize  \setcounter{utterance}{1}
\setlength{\tabcolsep}{0pt}
\begin{supertabular}{c@{$\;$}|p{.15\linewidth}@{}p{.15\linewidth}p{.15\linewidth}p{.15\linewidth}p{.15\linewidth}p{.15\linewidth}}
    \# & $\;$A & \multicolumn{4}{c}{Game Master} & $\;\:$B\\
    \hline

    \theutterance \stepcounter{utterance}  
    & & \multicolumn{4}{p{0.6\linewidth}}{
        \cellcolor[rgb]{0.9,0.9,0.9}{
            \makecell[{{p{\linewidth}}}]{
                \texttt{\tiny{[P1$\langle$GM]}}
                \texttt{Участвате в съвместна игра за преговори.} \\
\\ 
\texttt{Заедно с друг участник трябва да се споразумеете за един набор от предмети, които ще бъдат запазени. Всеки от вас има собствено мнение за това колко важен е за него всеки предмет (важност). Не знаете как другият участник оценява предметите. Освен това ви е дадено усилието, което изисква всеки предмет.} \\
\texttt{Можете да се споразумеете за набор, само ако общото усилие на избраните предмети не надвишава споделена граница:} \\
\\ 
\texttt{LIMIT = 5638} \\
\\ 
\texttt{Ето стойностите на усилието на отделните предмети:} \\
\\ 
\texttt{Усилие на предмета = \{"C08": 110, "C55": 727, "A43": 527, "A67": 855, "B42": 942, "A62": 617, "B19": 301, "C43": 130, "C73": 916, "C10": 212, "B36": 146, "A18": 559, "C61": 933, "C67": 740, "C42": 33, "C70": 799, "B40": 324, "C39": 841, "A21": 925, "C16": 639\}} \\
\\ 
\texttt{Ето вашето лично мнение за важността на всеки предмет:} \\
\\ 
\texttt{Стойности на важността на предметите = \{"C08": 204, "C55": 821, "A43": 621, "A67": 949, "B42": 1036, "A62": 711, "B19": 395, "C43": 224, "C73": 1010, "C10": 306, "B36": 240, "A18": 653, "C61": 1027, "C67": 834, "C42": 127, "C70": 893, "B40": 418, "C39": 935, "A21": 1019, "C16": 733\}} \\
\\ 
\texttt{Цел:} \\
\\ 
\texttt{Вашата цел е да договорите общ набор от предмети, който да ви е от възможно най{-}голяма полза (т.е. да максимизира общата важност за ВАС), като същевременно не надхвърляте LIMIT. Не сте длъжни да правите ПРЕДЛОЖЕНИЕ във всяко съобщение {-} можете и просто да преговаряте. Всички тактики са позволени!} \\
\\ 
\texttt{Протокол за взаимодействие:} \\
\\ 
\texttt{Можете да използвате само следните структурирани формати в едно съобщение:} \\
\\ 
\texttt{ПРЕДЛОЖЕНИЕ: \{'A', 'B', 'C', …\}} \\
\texttt{Предложете запазване точно на тези предмети.} \\
\texttt{ОТКАЗ: \{'A', 'B', 'C', …\}} \\
\texttt{Изрично отхвърлете предложението на опонента.} \\
\texttt{АРГУМЕНТ: \{'...'\}} \\
\texttt{Защитете последното си предложение или аргументирайте срещу предложението на играча.} \\
\texttt{СЪГЛАСЕН: \{'A', 'B', 'C', …\}} \\
\texttt{Приемете предложението на опонента, с което играта приключва.} \\
\\ 
\\ 
\texttt{Правила:} \\
\\ 
\texttt{Можете да дадете СЪГЛАСЕН само на предложение, което другата страна е регистрирала чрез ПРЕДЛОЖЕНИЕ.} \\
\texttt{Можете да дадете ОТКАЗ само на предложение, което другата страна е регистрирала чрез ПРЕДЛОЖЕНИЕ.} \\
\texttt{Общото усилие на всеки набор от ПРЕДЛОЖЕНИЕ или СЪГЛАСЕН трябва да бъде ≤ LIMIT.} \\
\texttt{НЕ разкривайте скритите си оценки за важност.} \\
\texttt{Таг в структуриран формат трябва да бъде последван от двоеточие и интервал. Аргументът трябва да бъде python набор, съдържащ 0 или повече низове.} \\
\texttt{Така че той трябва да бъде с формата "TAG: \{...\}"} \\
\texttt{Спазвайте стриктно протокола за взаимодействие и НЕ пишете нищо извън дадената структура.} \\
\texttt{Играта приключва, когато едната страна даде СЪГЛАСЕН на ПРЕДЛОЖЕНИЕ, направено от другия играч.} \\
\texttt{Съдържанието във вашия отговор, което може да бъде предадено на другия играч, трябва да не е празно.} \\
\texttt{Активни са само предложенията, които са били регистрирани чрез структурата на формата ПРЕДЛОЖЕНИЕ и които не са били отхвърлени чрез ОТКАЗ.} \\
\\ 
            }
        }
    }
    & \\ \\

    \theutterance \stepcounter{utterance}  
    & \multicolumn{4}{p{0.6\linewidth}}{
        \cellcolor[rgb]{0.9,0.9,0.9}{
            \makecell[{{p{\linewidth}}}]{
                \texttt{\tiny{[P1$\rangle$GM]}}
                \texttt{ПРЕДЛОЖЕНИЕ: \{'C08', 'A43', 'B19', 'C43', 'C10', 'B36', 'C42', 'B40', 'C16'\}} \\
            }
        }
    }
    & & \\ \\

    \theutterance \stepcounter{utterance}  
    & & \multicolumn{4}{p{0.6\linewidth}}{
        \cellcolor[rgb]{0.9,0.9,0.9}{
            \makecell[{{p{\linewidth}}}]{
                \texttt{\tiny{[GM$\rangle$P2]}}
                \texttt{Участвате в съвместна игра за преговори.} \\
\\ 
\texttt{Заедно с друг участник трябва да се споразумеете за един набор от предмети, които ще бъдат запазени. Всеки от вас има собствено мнение за това колко важен е за него всеки предмет (важност). Не знаете как другият участник оценява предметите. Освен това ви е дадено усилието, което изисква всеки предмет.} \\
\texttt{Можете да се споразумеете за набор, само ако общото усилие на избраните предмети не надвишава споделена граница:} \\
\\ 
\texttt{LIMIT = 5638} \\
\\ 
\texttt{Ето стойностите на усилието на отделните предмети:} \\
\\ 
\texttt{Усилие на предмета = \{"C08": 110, "C55": 727, "A43": 527, "A67": 855, "B42": 942, "A62": 617, "B19": 301, "C43": 130, "C73": 916, "C10": 212, "B36": 146, "A18": 559, "C61": 933, "C67": 740, "C42": 33, "C70": 799, "B40": 324, "C39": 841, "A21": 925, "C16": 639\}} \\
\\ 
\texttt{Тук е вашето лично мнение за важността на всеки предмет:} \\
\\ 
\texttt{Стойности на важността на предметите = \{"C08": 204, "C55": 821, "A43": 621, "A67": 949, "B42": 1036, "A62": 711, "B19": 395, "C43": 224, "C73": 1010, "C10": 306, "B36": 240, "A18": 653, "C61": 1027, "C67": 834, "C42": 127, "C70": 893, "B40": 418, "C39": 935, "A21": 1019, "C16": 733\}} \\
\\ 
\texttt{Цел:} \\
\\ 
\texttt{Вашата цел е да договорите общ набор от предмети, който да ви е от възможно най{-}голяма полза (т.е. да максимизира общата важност за ВАС), като същевременно не надхвърляте LIMIT. Не сте длъжни да правите ПРЕДЛОЖЕНИЕ във всяко съобщение {-} можете и просто да преговаряте. Всички тактики са позволени!} \\
\\ 
\texttt{Протокол за взаимодействие:} \\
\\ 
\texttt{Можете да използвате само следните структурирани формати в едно съобщение:} \\
\\ 
\texttt{ПРЕДЛОЖЕНИЕ: \{'A', 'B', 'C', …\}} \\
\texttt{Предложете запазване точно на тези предмети.} \\
\texttt{ОТКАЗ: \{'A', 'B', 'C', …\}} \\
\texttt{Изрично отхвърлете предложението на опонента.} \\
\texttt{АРГУМЕНТ: \{'...'\}} \\
\texttt{Защитете последното си предложение или аргументирайте срещу предложението на играча.} \\
\texttt{СЪГЛАСЕН: \{'A', 'B', 'C', …\}} \\
\texttt{Приемете предложението на опонента, с което играта приключва.} \\
\\ 
\\ 
\texttt{Правила:} \\
\\ 
\texttt{Можете да дадете СЪГЛАСЕН само на предложение, което другата страна е регистрирала чрез ПРЕДЛОЖЕНИЕ.} \\
\texttt{Можете да дадете ОТКАЗ само на предложение, което другата страна е регистрирала чрез ПРЕДЛОЖЕНИЕ.} \\
\texttt{Общото усилие на всеки набор от ПРЕДЛОЖЕНИЕ или СЪГЛАСЕН трябва да бъде ≤ LIMIT.} \\
\texttt{НЕ разкривайте скритите си оценки за важност.} \\
\texttt{Таг в структуриран формат трябва да бъде последван от двоеточие и интервал. Аргументът трябва да бъде python набор, съдържащ 0 или повече низове.} \\
\texttt{Така че той трябва да бъде с формата "TAG: \{...\}"} \\
\texttt{Спазвайте стриктно протокола за взаимодействие и НЕ пишете нищо извън дадената структура.} \\
\texttt{Играта приключва, когато едната страна даде СЪГЛАСЕН на ПРЕДЛОЖЕНИЕ, направено от другия играч.} \\
\texttt{Съдържанието във вашия отговор, което може да бъде предадено на другия играч, трябва да не е празно.} \\
\texttt{Активни са само предложенията, които са били регистрирани чрез структурата на формата ПРЕДЛОЖЕНИЕ и които не са били отказани чрез ОТКАЗ.} \\
\\ 
\\ 
\\ 
\texttt{Сега ще получите първото съобщение на другия играч.} \\
\\ 
\texttt{ПРЕДЛОЖЕНИЕ: \{'C08', 'A43', 'B19', 'C43', 'C10', 'B36', 'C42', 'B40', 'C16'\}} \\
            }
        }
    }
    & \\ \\

    \theutterance \stepcounter{utterance}  
    & & & \multicolumn{4}{p{0.6\linewidth}}{
        \cellcolor[rgb]{0.9,0.9,0.9}{
            \makecell[{{p{\linewidth}}}]{
                \texttt{\tiny{[GM$\langle$P2]}}
                \texttt{ОТКАЗ: \{'C08', 'A43', 'B19', 'C43', 'C10', 'B36', 'C42', 'B40', 'C16'\}} \\
\\ 
\texttt{АРГУМЕНТ: \{'Това предложение включва предмети с ниска важност за мен, като C08, C43 и C42, които имат относително високо усилие спрямо ползата им. Предлагам да се фокусираме върху предмети с по{-}висока важност и по{-}ефективно използване на усилието.'\}} \\
            }
        }
    }
     \\ \\

    \theutterance \stepcounter{utterance}  
    & & \multicolumn{4}{p{0.6\linewidth}}{
        \cellcolor[rgb]{0.9,0.9,0.9}{
            \makecell[{{p{\linewidth}}}]{
                \texttt{\tiny{[P1$\langle$GM]}}
                \texttt{ОТКАЗ: \{'C08', 'A43', 'B19', 'C43', 'C10', 'B36', 'C42', 'B40', 'C16'\}} \\
\\ 
\texttt{АРГУМЕНТ: \{'Това предложение включва предмети с ниска важност за мен, като C08, C43 и C42, които имат относително високо усилие спрямо ползата им. Предлагам да се фокусираме върху предмети с по{-}висока важност и по{-}ефективно използване на усилието.'\}} \\
            }
        }
    }
    & \\ \\

    \theutterance \stepcounter{utterance}  
    & \multicolumn{4}{p{0.6\linewidth}}{
        \cellcolor[rgb]{0.9,0.9,0.9}{
            \makecell[{{p{\linewidth}}}]{
                \texttt{\tiny{[P1$\rangle$GM]}}
                \texttt{ПРЕДЛОЖЕНИЕ: \{'C55', 'A43', 'A62', 'B19', 'C10', 'B36', 'A18', 'C16'\}} \\
\\ 
\texttt{АРГУМЕНТ: \{'Това предложение включва предмети с висока важност за мен (като C55, A43, A62 и A18), като същевременно запазва баланс по усилие. Общото усилие е 3916, което е значително под LIMIT и оставя пространство за допълнителни предмети, ако желаете.'\}} \\
            }
        }
    }
    & & \\ \\

    \theutterance \stepcounter{utterance}  
    & & \multicolumn{4}{p{0.6\linewidth}}{
        \cellcolor[rgb]{0.9,0.9,0.9}{
            \makecell[{{p{\linewidth}}}]{
                \texttt{\tiny{[GM$\rangle$P2]}}
                \texttt{ПРЕДЛОЖЕНИЕ: \{'C55', 'A43', 'A62', 'B19', 'C10', 'B36', 'A18', 'C16'\}} \\
\\ 
\texttt{АРГУМЕНТ: \{'Това предложение включва предмети с висока важност за мен (като C55, A43, A62 и A18), като същевременно запазва баланс по усилие. Общото усилие е 3916, което е значително под LIMIT и оставя пространство за допълнителни предмети, ако желаете.'\}} \\
            }
        }
    }
    & \\ \\

    \theutterance \stepcounter{utterance}  
    & & & \multicolumn{4}{p{0.6\linewidth}}{
        \cellcolor[rgb]{0.9,0.9,0.9}{
            \makecell[{{p{\linewidth}}}]{
                \texttt{\tiny{[GM$\langle$P2]}}
                \texttt{ПРЕДЛОЖЕНИЕ: \{'C55', 'A43', 'A67', 'B42', 'A62', 'B19', 'C10', 'B36', 'A18'\}} \\
\\ 
\texttt{АРГУМЕНТ: \{'Включих A67 и B42, които са с висока важност за мен, като същевременно запазих баланса по усилие (общо 5160). Това предложение подобрява значително ползата за мен, без да надвишава LIMIT.'\}} \\
            }
        }
    }
     \\ \\

    \theutterance \stepcounter{utterance}  
    & & \multicolumn{4}{p{0.6\linewidth}}{
        \cellcolor[rgb]{0.9,0.9,0.9}{
            \makecell[{{p{\linewidth}}}]{
                \texttt{\tiny{[P1$\langle$GM]}}
                \texttt{ПРЕДЛОЖЕНИЕ: \{'C55', 'A43', 'A67', 'B42', 'A62', 'B19', 'C10', 'B36', 'A18'\}} \\
\\ 
\texttt{АРГУМЕНТ: \{'Включих A67 и B42, които са с висока важност за мен, като същевременно запазих баланса по усилие (общо 5160). Това предложение подобрява значително ползата за мен, без да надвишава LIMIT.'\}} \\
            }
        }
    }
    & \\ \\

    \theutterance \stepcounter{utterance}  
    & \multicolumn{4}{p{0.6\linewidth}}{
        \cellcolor[rgb]{0.9,0.9,0.9}{
            \makecell[{{p{\linewidth}}}]{
                \texttt{\tiny{[P1$\rangle$GM]}}
                \texttt{ОТКАЗ: \{'C55', 'A43', 'A67', 'B42', 'A62', 'B19', 'C10', 'B36', 'A18'\}} \\
\\ 
\texttt{АРГУМЕНТ: \{'Предложението ви включва предмети с много високо усилие (като B42 и A67), които ограничават възможността за добавяне на други ценни за мен предмети. Предлагам алтернатива, която запазва по{-}висока гъвкавост и по{-}добър баланс между важност и усилие.'\}} \\
\\ 
\texttt{ПРЕДЛОЖЕНИЕ: \{'C55', 'A43', 'A62', 'B19', 'C73', 'C10', 'B36', 'A18', 'C16'\}} \\
\texttt{(Общо усилие: 5060)} \\
            }
        }
    }
    & & \\ \\

    \theutterance \stepcounter{utterance}  
    & & \multicolumn{4}{p{0.6\linewidth}}{
        \cellcolor[rgb]{0.9,0.9,0.9}{
            \makecell[{{p{\linewidth}}}]{
                \texttt{\tiny{[P1$\langle$GM]}}
                \texttt{Използвахте таг ПРЕДЛОЖЕНИЕ, но не предоставихте валиден python set, съдържащ низове като аргументи, например \{'A', 'B', 'C', ...\}. Опитайте отново.} \\
            }
        }
    }
    & \\ \\

    \theutterance \stepcounter{utterance}  
    & \multicolumn{4}{p{0.6\linewidth}}{
        \cellcolor[rgb]{0.9,0.9,0.9}{
            \makecell[{{p{\linewidth}}}]{
                \texttt{\tiny{[P1$\rangle$GM]}}
                \texttt{ПРЕДЛОЖЕНИЕ: \{'C55', 'A43', 'A62', 'B19', 'C73', 'C10', 'B36', 'A18', 'C16'\}} \\
\\ 
\texttt{АРГУМЕНТ: \{'Това предложение включва C73 (висока важност за мен) и запазва баланса по усилие (5060). То е по{-}ефективно от последното ви предложение, тъй като оставя повече пространство за допълнителни предмети или компромиси.'\}} \\
            }
        }
    }
    & & \\ \\

    \theutterance \stepcounter{utterance}  
    & & \multicolumn{4}{p{0.6\linewidth}}{
        \cellcolor[rgb]{0.9,0.9,0.9}{
            \makecell[{{p{\linewidth}}}]{
                \texttt{\tiny{[GM$\rangle$P2]}}
                \texttt{ПРЕДЛОЖЕНИЕ: \{'C55', 'A43', 'A62', 'B19', 'C73', 'C10', 'B36', 'A18', 'C16'\}} \\
\\ 
\texttt{АРГУМЕНТ: \{'Това предложение включва C73 (висока важност за мен) и запазва баланса по усилие (5060). То е по{-}ефективно от последното ви предложение, тъй като оставя повече пространство за допълнителни предмети или компромиси.'\}} \\
            }
        }
    }
    & \\ \\

    \theutterance \stepcounter{utterance}  
    & & & \multicolumn{4}{p{0.6\linewidth}}{
        \cellcolor[rgb]{0.9,0.9,0.9}{
            \makecell[{{p{\linewidth}}}]{
                \texttt{\tiny{[GM$\langle$P2]}}
                \texttt{ОТКАЗ: \{'C55', 'A43', 'A62', 'B19', 'C73', 'C10', 'B36', 'A18', 'C16'\}} \\
\\ 
\texttt{АРГУМЕНТ: \{'C73 има твърде високо усилие (916) за относителната си важност в моя контекст. Предпочитам да запазим предмети с по{-}добро съотношение важност/усилие. Предлагам да заменим C73 с A67 (855 усилие, но с по{-}висока важност за мен).'\}} \\
\\ 
\texttt{ПРЕДЛОЖЕНИЕ: \{'C55', 'A43', 'A67', 'B42', 'A62', 'B19', 'C10', 'B36', 'A18'\}} \\
\texttt{Общо усилие: **5159** (под LIMIT)} \\
            }
        }
    }
     \\ \\

    \theutterance \stepcounter{utterance}  
    & & \multicolumn{4}{p{0.6\linewidth}}{
        \cellcolor[rgb]{0.9,0.9,0.9}{
            \makecell[{{p{\linewidth}}}]{
                \texttt{\tiny{[GM$\rangle$P2]}}
                \texttt{Използвахте таг ПРЕДЛОЖЕНИЕ, но не предоставихте валиден python set, съдържащ низове като аргументи, например \{'A', 'B', 'C', ...\}. Опитайте отново.} \\
            }
        }
    }
    & \\ \\

    \theutterance \stepcounter{utterance}  
    & & & \multicolumn{4}{p{0.6\linewidth}}{
        \cellcolor[rgb]{0.9,0.9,0.9}{
            \makecell[{{p{\linewidth}}}]{
                \texttt{\tiny{[GM$\langle$P2]}}
                \texttt{ПРЕДЛОЖЕНИЕ: \{'C55', 'A43', 'A67', 'B42', 'A62', 'B19', 'C10', 'B36', 'A18'\}} \\
\\ 
\texttt{АРГУМЕНТ: \{'Общото усилие на това предложение е 5159, което е под LIMIT. То включва предмети с висока важност за мен (A67, B42) и запазва баланс по усилие.'\}} \\
            }
        }
    }
     \\ \\

    \theutterance \stepcounter{utterance}  
    & & \multicolumn{4}{p{0.6\linewidth}}{
        \cellcolor[rgb]{0.9,0.9,0.9}{
            \makecell[{{p{\linewidth}}}]{
                \texttt{\tiny{[P1$\langle$GM]}}
                \texttt{ПРЕДЛОЖЕНИЕ: \{'C55', 'A43', 'A67', 'B42', 'A62', 'B19', 'C10', 'B36', 'A18'\}} \\
\\ 
\texttt{АРГУМЕНТ: \{'Общото усилие на това предложение е 5159, което е под LIMIT. То включва предмети с висока важност за мен (A67, B42) и запазва баланс по усилие.'\}} \\
            }
        }
    }
    & \\ \\

    \theutterance \stepcounter{utterance}  
    & \multicolumn{4}{p{0.6\linewidth}}{
        \cellcolor[rgb]{0.9,0.9,0.9}{
            \makecell[{{p{\linewidth}}}]{
                \texttt{\tiny{[P1$\rangle$GM]}}
                \texttt{ОТКАЗ: \{'C55', 'A43', 'A67', 'B42', 'A62', 'B19', 'C10', 'B36', 'A18'\}} \\
\\ 
\texttt{АРГУМЕНТ: \{'Въпреки че усилието е под LIMIT, включването на A67 (855) и B42 (942) заема твърде много от ограничението, което ограничава възможността за добавяне на други ценни за мен предмети като C73 или C61. Предлагам да заменим A67 и B42 с по{-}ефективни алтернативи.'\}} \\
\\ 
\texttt{ПРЕДЛОЖЕНИЕ: \{'C55', 'A43', 'A62', 'B19', 'C73', 'C10', 'B36', 'A18', 'C67'\}} \\
\texttt{(Общо усилие: 5088)} \\
            }
        }
    }
    & & \\ \\

    \theutterance \stepcounter{utterance}  
    & & & \multicolumn{2}{p{0.3\linewidth}}{
        \cellcolor[rgb]{0.9,0.9,0.9}{
            \makecell[{{p{\linewidth}}}]{
                \texttt{\tiny{[GM$|$GM]}}
                \texttt{abort} \\
            }
        }
    }
    & & \\ \\

    \theutterance \stepcounter{utterance}  
    & & & \multicolumn{2}{p{0.3\linewidth}}{
        \cellcolor[rgb]{0.9,0.9,0.9}{
            \makecell[{{p{\linewidth}}}]{
                \texttt{\tiny{[GM$|$GM]}}
                \texttt{end game} \\
            }
        }
    }
    & & \\ \\

\end{supertabular}
}

\end{document}
